\chapter{SOLUCIÓN IMPLEMENTADA}
\label{ch:solucion}

Este capítulo describe, de forma general, la solución diseñada antes de entrar
al detalle teórico y de implementación. A partir de las alternativas revisadas
en el Capítulo~\ref{ch:estado_arte}, se propone un robot móvil diferencial de
bajo costo que integra percepción láser 2D, odometría por encoders y control
PID embebido, cuyas ganancias se determinaron combinando optimización por
enjambre de partículas (PSO) y refinamiento experimental sobre el robot, sobre
una arquitectura distribuida en dos entornos: robot y
estación de trabajo, comunicados mediante \gls{ROS2}. El objetivo es que el
sistema sea capaz de mapear, navegar y explorar entornos interiores de forma
autónoma.
La Fig.~\ref{fig:arquitectura_alto_nivel} presenta el diagrama de bloques
simplificado del sistema, organizado en cuatro dominios funcionales
identificados por color.

\utemFigurePropia{2.Texto/Imag/arquitectura_alto_nivel.png}{0.94\textwidth}
{Diagrama de bloques simplificado de la solución implementada.}
{fig:arquitectura_alto_nivel}

Las flechas del diagrama representan el flujo lógico de mediciones, estados,
metas y comandos entre los componentes. Por tanto, no todas corresponden a
conexiones físicas directas, ya que parte del intercambio se realiza mediante
tópicos, acciones y transformaciones de \gls{ROS2} a través de la red local.
La función general de cada dominio es la siguiente:

\begin{description}
\item[Sensado.] Los \emph{encoders} de cuadratura miden el giro de cada rueda
y el \gls{LIDAR} 2D percibe el entorno; son las dos fuentes de información del
sistema.
\item[Control de bajo nivel.] El \gls{Arduino} lee los encoders y ejecuta el
control PID de velocidad de cada rueda, accionando los motores DC de tracción
diferencial. Las ganancias de ese PID se determinaron en dos etapas: PSO
entregó un primer conjunto a partir de un modelo identificado de cada rueda, y
ese conjunto se refinó luego sobre el robot real hasta llegar a los valores
definitivos.
\item[Sistema computacional del robot.] La Raspberry Pi~4 integra \gls{ROS2},
estima la odometría a partir de las velocidades de rueda y actúa de puente
entre el hardware embebido y la estación de trabajo.
\item[Mapeo, navegación y exploración.] En la estación de trabajo,
\texttt{slam\_toolbox} construye el mapa y estima la pose, Nav2 planifica y
controla la navegación, el explorador de fronteras genera metas autónomas y
RViz2 permite la supervisión.
\end{description}

\section{Selección de tecnologías}
\label{sec:seleccion_tecnologias}

Cada dominio del sistema se resolvió eligiendo, entre las alternativas
comparadas en el Capítulo~\ref{ch:estado_arte}, la más adecuada para una
plataforma diferencial de bajo costo orientada a interiores. Se adoptó
\gls{ROS2} como middleware por su ecosistema de paquetes de mapeo, navegación
y exploración y por su soporte vigente para sistemas distribuidos
(Sec.~\ref{sec:alternativas_middleware})
\citep{macenski2022ros2,ros_noetic_eol}. La percepción se resolvió con un
\gls{LIDAR} 2D, suficiente para interiores y de menor costo y carga de
procesamiento que las cámaras RGB-D o el \gls{LIDAR} 3D, y el mapeo con
\texttt{slam\_toolbox}, por su integración nativa, su enfoque basado en grafos
de poses y su capacidad de guardar y cargar mapas
(Sec.~\ref{sec:sensores}) \citep{macenski2021slamtoolbox}. La navegación se
apoyó en Nav2, con un planificador global de búsqueda sobre rejilla y un
controlador local reactivo, adecuados a un mapa de ocupación 2D y de bajo costo
computacional (Sec.~\ref{sec:exploracion})
\citep{hart1968astar,macenski2023desks}.

 La validación previa del diseño se apoyó en un gemelo digital construido
en Gazebo Classic. Se optó por esa versión y no por el Gazebo actual porque
es la compatible con \gls{ROS2} Foxy y con Ubuntu 20.04, que constituyen la
base del sistema (Sec.~\ref{sec:gemelo_digital}).

La sintonización del control de bajo nivel también se resolvió como una
decisión de diseño. En lugar de buscar las tres ganancias PID por prueba y
error desde cero, se identificó un modelo de primer orden con retardo para cada
rueda y sobre él
se aplicó PSO, un método de optimización que no requiere derivadas ni un
modelo exacto de la planta y que permite buscar las tres ganancias de forma
simultánea. El resultado de esa búsqueda se adoptó como punto de partida y no
como valor final: las ganancias definitivas se alcanzaron ajustando ese punto
de partida sobre el robot real, donde intervienen efectos que el modelo no
recoge, como la zona muerta del driver, la cuantización de los encoders y la
asimetría entre ambas ruedas.

Para la exploración autónoma se eligió el enfoque basado en fronteras
(Sec.~\ref{sec:exploracion}), por su menor complejidad y su fácil integración
con el mapa de ocupación y con Nav2, frente a los métodos de ganancia de
información y de \textit{active SLAM}, que exigen modelos y cómputo adicionales
\citep{liu2025frontier,horner_mexplore}.

El fundamento teórico de cada uno de estos componentes se desarrolla en el
Capítulo~\ref{marco_teorico}, y su implementación concreta en el
Capítulo~\ref{ch:implementacion}.